\documentclass[12pt]{article}

\usepackage[utf8]{inputenc} % important for pdflatex + any UTF-8 characters
\usepackage{amsmath, amssymb}
\usepackage{geometry}
\geometry{margin=1in}

\setlength{\parindent}{0pt}

\title{Solutions -- Practice Problem Set 2}
\author{ECON 300 -- Labour Economics}
\date{Winter 2026}

\begin{document}
\maketitle


\subsection*{Problem 1: }

\textbf{Program parameters:}
\[
G = 600,\quad t = 0.5,\quad w = 20.
\]
While the individual is on the benefit, disposable income is:
\[
C = G + (1-t)wh \quad \text{(as long as benefits are positive)}.
\]

\begin{enumerate}
\item \textbf{Break-even level of earnings.}

Benefits fall to zero when the reduction in benefits equals the guarantee. With a tax-back rate $t$, benefits are reduced by $t \cdot (wh)$ when earnings are $wh$.
Break-even occurs when:
\[
t(wh) = G.
\]
Solve for earnings:
\[
wh = \frac{G}{t} = \frac{600}{0.5} = 1200.
\]
So the break-even level of \textbf{earnings} is \$1200 per week.

If you want the break-even \textbf{hours}, divide by the wage:
\[
h^{BE} = \frac{1200}{20} = 60.
\]
So benefits reach zero at \textbf{60 hours} per week.

\item \textbf{Effective net wage while receiving benefits.}

When benefits are being reduced at rate $t$, each additional dollar of earnings reduces benefits by $t$ dollars.
So the net gain from an extra dollar of earnings is $(1-t)$.

Therefore, the \textbf{effective net wage} is:
\[
w_{\text{net}} = (1-t)w = (1-0.5)\cdot 20 = 10.
\]

\item \textbf{Comparison and implication for labour supply incentives.}

The market wage is \$20, but while receiving benefits the worker effectively keeps only \$10 per hour.
That means the program reduces the slope of the budget constraint during the benefit range and lowers the opportunity cost of leisure (relative to the no-program case).
Mechanically, a lower net wage reduces the substitution incentive to work and tends to reduce hours of work (especially on the intensive margin).

\item \textbf{Why high benefit reduction rates may discourage work.}

A high tax-back rate $t$ implies a low net-of-tax share $(1-t)$.
As $t$ rises toward 1, the effective net wage $(1-t)w$ falls toward zero.
This makes additional work yield little increase in disposable income, weakening work incentives.
In extreme cases, if $t=1$ then additional earnings reduce benefits dollar-for-dollar, so disposable income does not rise with work, creating strong disincentives to increase hours.
\end{enumerate}

\subsection*{Problem 2: }

Let total available time be $T=80$ hours per week. Let leisure be $L$ and hours of work be:
\[
h = T - L.
\]

\textbf{(i) No income support.} If there is no non-labour income, the budget constraint is:
\[
C = w(T-L).
\]
In $(L,C)$ space this is a straight line with:
\begin{itemize}
\item Leisure intercept at $L=T$ (if $h=0$, then $L=T$),
\item Consumption intercept at $C=wT$ (if $L=0$, work all time),
\item Slope $-w$.
\end{itemize}

\textbf{(ii) Guaranteed income with tax-back rate $t<1$.}
A standard income maintenance program with guarantee $G$ and reduction rate $t$ generates a \emph{kinked} budget constraint:
\[
C =
\begin{cases}
G + (1-t)w(T-L), & \text{if benefits are positive},\\
w(T-L), & \text{after benefits are exhausted}.
\end{cases}
\]
Key features:
\begin{itemize}
\item At $h=0$ (i.e., $L=T$), consumption equals $G$.
\item While on benefit, slope is $-(1-t)w$, which is flatter than $-w$ because $(1-t)w<w$.
\item After the break-even point, the slope returns to $-w$ (no benefits).
\end{itemize}

\begin{enumerate}
\item \textbf{Draw both constraints.} 
Plot $C=w(T-L)$ for the no-program case. For the program case, start at $(L=T,C=G)$ and draw a line with slope $-(1-t)w$ until the break-even point, then continue with slope $-w$.

\item \textbf{Label $G$ and the break-even point.}
\begin{itemize}
\item The guaranteed income $G$ is the vertical intercept of the program constraint at $L=T$.
\item Break-even occurs where benefit reaches zero, i.e., $twh=G$ (equivalently $t w(T-L)=G$).
\end{itemize}

\item \textbf{Slope changes.} The slope of the budget constraint is the \emph{net wage} in consumption units per hour of leisure.
\begin{itemize}
\item No-program slope: $-w$.
\item Program slope while on benefit: $-(1-t)w$ (flatter).
\item After break-even: back to $-w$.
\end{itemize}
So, the program reduces the net return to working within the benefit region.

\item \textbf{Indifference curves and predicted labour supply effects.}

Two margins matter:
\begin{itemize}
\item \textbf{Extensive margin (participation):} Because consumption at $h=0$ rises from $0$ (or from any baseline) to $G$, some individuals who would otherwise work might choose not to participate, especially if their preferred choice is near the endowment point.
\item \textbf{Intensive margin (hours):} For those who still work, the flatter slope $(1-t)w$ reduces the substitution incentive to work. This tends to lower hours while on benefit. However, the income effect of higher resources can also increase leisure further if leisure is normal.
\end{itemize}
Overall, income support typically reduces labour supply, but the magnitude depends on preferences and the program design.
\end{enumerate}

\subsection*{Problem 3: }

Answer each briefly (3--5 sentences).

\begin{enumerate}
\item \textbf{Negative income tax (NIT) vs traditional welfare.}

A negative income tax is usually framed as a unified schedule where individuals receive a transfer when income is below a threshold, with benefits reduced smoothly as earnings rise.
Traditional welfare programs often have categorical eligibility, complex rules, and may feature sharp cutoffs that can create large effective marginal tax rates.
Conceptually, the NIT emphasizes a simpler, more continuous budget constraint, reducing ``welfare wall'' effects relative to abrupt benefit loss.

\item \textbf{Why an EITC may increase participation but reduce hours for some.}

An earned income tax credit typically \emph{subsidizes} earnings in the phase-in region, raising the effective wage and encouraging labour force entry (higher participation).
In the phase-out region, the credit is withdrawn as earnings rise, which raises effective marginal tax rates and can reduce hours for some workers who are already employed.
Thus, EITC designs often increase labour supply on the extensive margin but can have ambiguous or negative effects on hours for higher earners within the eligible range.

\item \textbf{Extensive vs intensive margin.}

The extensive margin refers to the decision of whether to work at all (participation), while the intensive margin refers to how many hours to work conditional on participating.
Policies like fixed benefits or entry bonuses mainly affect participation, whereas changes in net wages (tax rates, benefit reduction rates) affect hours.
Empirically and theoretically, these margins can respond differently, so separating them helps interpret the total labour supply response.

\item \textbf{Equity-efficiency trade-offs with targeting.}

Targeting benefits to low-income workers improves equity by directing resources to those with the greatest need and often reduces program cost.
However, targeting typically requires benefits to be reduced as earnings rise, which increases effective marginal tax rates and can reduce work incentives (efficiency cost).
Hence, better redistribution often comes at the cost of weaker labour supply incentives unless the program is designed to mitigate high implicit taxes.
\end{enumerate}



\subsection*{Problem 5: }

\textbf{Production function:}
\[
Q = 10L^{0.5} = 10\sqrt{L}.
\]
\textbf{Given:} $P=4$, $w=8$.

\begin{enumerate}
\item \textbf{Marginal product of labour (MPL).}

Compute the derivative:
\[
MPL = \frac{dQ}{dL} = \frac{d}{dL}\left(10L^{1/2}\right)
= 10 \cdot \frac{1}{2}L^{-1/2}
= 5L^{-1/2}
= \frac{5}{\sqrt{L}}.
\]

\item \textbf{Value of marginal product (VMPL).}

\[
VMPL = P \cdot MPL = 4 \cdot \frac{5}{\sqrt{L}} = \frac{20}{\sqrt{L}}.
\]

\item \textbf{Optimal employment level.}

In a competitive labour market, the firm hires labour until:
\[
VMPL = w.
\]
So:
\[
\frac{20}{\sqrt{L}} = 8
\quad \Rightarrow \quad
20 = 8\sqrt{L}
\quad \Rightarrow \quad
\sqrt{L} = \frac{20}{8} = 2.5
\quad \Rightarrow \quad
L^* = (2.5)^2 = 6.25.
\]

So the optimal employment is:
\[
L^* = 6.25.
\]
(Interpreting $L$ as a continuous measure of labour input; if $L$ must be integer workers, the firm would compare profits at $L=6$ and $L=7$.)

\item \textbf{Why labour demand is downward sloping in $w$.}

From the condition $VMPL=w$, we can solve for $L$ as a function of $w$:
\[
\frac{20}{\sqrt{L}} = w
\quad \Rightarrow \quad
\sqrt{L} = \frac{20}{w}
\quad \Rightarrow \quad
L(w) = \left(\frac{20}{w}\right)^2.
\]
As $w$ rises, $\frac{20}{w}$ falls, so $L(w)$ falls. This gives a downward-sloping labour demand curve.
Economically, because of diminishing marginal product, higher wages make additional labour less profitable, so the firm reduces employment.
\end{enumerate}

\subsection*{Problem 6: }

From Problem 5, the labour demand curve (for the firm) is given by:
\[
w = VMPL = \frac{20}{\sqrt{L}}
\quad \text{or equivalently} \quad
L(w) = \left(\frac{20}{w}\right)^2.
\]

\begin{enumerate}
\item \textbf{Draw the labour demand curve.}
In $(L,w)$ space, plot:
\[
w(L)=\frac{20}{\sqrt{L}},
\]
which is downward sloping and convex (flattens as $L$ increases).

\item \textbf{Equilibrium employment at $w=8$ and $w=12$.}

At $w=8$:
\[
L(8)=\left(\frac{20}{8}\right)^2 = (2.5)^2 = 6.25.
\]

At $w=12$:
\[
L(12)=\left(\frac{20}{12}\right)^2 = \left(\frac{5}{3}\right)^2 = \frac{25}{9} \approx 2.78.
\]

So when the wage rises from 8 to 12, employment falls from $6.25$ to about $2.78$.

\item \textbf{Wage change vs output price or technology change.}

A \textbf{wage change} causes a movement along a given labour demand curve because the curve itself is defined by $VMPL(L)$.

A change in \textbf{output price} $P$ changes VMPL at every $L$ because $VMPL=P\cdot MPL$. Higher $P$ increases VMPL and shifts labour demand outward (upward in $(L,w)$ space).

A change in \textbf{technology} changes the production function and MPL at every $L$. If technology raises productivity (higher MPL), VMPL rises and labour demand shifts outward as well.
\end{enumerate}

\subsection*{Problem 7: }

Now suppose output price increases from $P=4$ to $P=6$, holding the production function fixed.

\begin{enumerate}
\item \textbf{Recalculate optimal $L$ when $P=6$.}

We already have:
\[
MPL = \frac{5}{\sqrt{L}}.
\]
So new VMPL is:
\[
VMPL = P \cdot MPL = 6 \cdot \frac{5}{\sqrt{L}} = \frac{30}{\sqrt{L}}.
\]
Set $VMPL=w$ with $w=8$:
\[
\frac{30}{\sqrt{L}} = 8
\quad \Rightarrow \quad
30 = 8\sqrt{L}
\quad \Rightarrow \quad
\sqrt{L} = \frac{30}{8} = 3.75
\quad \Rightarrow \quad
L^* = (3.75)^2 = 14.0625.
\]

So:
\[
L^* \approx 14.06.
\]

\item \textbf{Intuition.}

When the output price rises, each unit of output is worth more, so the revenue generated by an additional unit of labour rises.
Because $VMPL=P\cdot MPL$, a higher $P$ increases VMPL at every $L$, making it profitable to hire more labour until the higher VMPL is driven back down to the wage through diminishing marginal product.

\item \textbf{Movement along vs shift.}

\begin{itemize}
\item A \textbf{movement along} labour demand occurs when $w$ changes holding $P$ and technology fixed. The firm changes $L$ because the hiring condition $VMPL=w$ changes due to $w$.
\item A \textbf{shift} in labour demand occurs when $P$ or technology changes, because $VMPL(L)$ changes at every $L$. Here, higher $P$ shifts labour demand outward (upward): at any given $L$, the firm is willing to pay a higher wage, and at any given wage, it hires more labour.
\end{itemize}
\end{enumerate}

\vfill
\textbf{End of Solutions}

\end{document}
